\section{A sparse estimate for packet arrays}
\label{sec:sparse}

The transverse saving has a useful consequence for arrays whose individual
packets see little measure. The broad--narrow decomposition follows the
framework of Bourgain--Guth~\cite[Section~2]{BG11} and its fractal
restriction implementation in Du--Zhang~\cite[Section~3]{DZ19}; the
parabolic rescaling below also preserves packet sparsity. Broad
contributions gain a power by
Proposition~\ref{prop:transverse}; when the iteration stops, the sparsity
condition supplies the gain through the refined fractal estimate.

\begin{proposition}[Sparse packet estimate]\label{prop:sparse}
There is an absolute constant $c_s>0$ with the following property.  Fix
$\tau_0>0$, and let $L^{\tau_0}\le D\le L$.  Let $U$ be a finite packet
array of the form \eqref{eq:src5}, with coefficient vector $c$, and let
$\mu$ be a nonnegative measure such that
\[
 \mu(B(z,r))\le m r^2\quad(r\ge1),\qquad
 \supp\mu\subset\{|t|\le C L^2\}.
\]
Assume the packet sparsity condition \eqref{eq:src6}, with $M=100$, for
every label with nonzero coefficient.  Then
\begin{equation}\label{eq:src7}
 \sum_{q\text{ a unit lattice cube}}
       \mu(q)^{2/3}\sup_q|U|^2
 \lesssim_{\tau_0}
       m^{2/3}L^{4/3}D^{-c_s}\|c\|_{\ell^2}^2.
\end{equation}
Only finitely many uniform profile derivative bounds, depending on
$\tau_0$, are required.  Consequently the same assertion holds for any
prescribed finite set of field derivatives, provided the original
profiles have sufficiently many bounded derivatives.  Fixed enlargements
of the time range and restriction to a submeasure are allowed.  The
exponent $c_s$ is independent of $\tau_0$; the implicit constant and the
required differentiability need not be.
\end{proposition}

The proof is divided into the exact-field reduction, a broad estimate, and
the rescaling iteration.  All bounds on packet classes and frequency
boxes are fixed throughout.  Zero measures and zero coefficient vectors
are immediate, so normalization by their sizes is legitimate.

\subsection{A local norm and an exact-field reduction}

Choose a dyadic integer $K\asymp L^\beta$, with $\beta>0$ to be specified
at the end.  For a measure $\nu$ define
\begin{equation}\label{eq:sparse-functional}
 \mathcal T_\nu(V)
   =\sum_{Q\text{ a }K^2\text{-cube}}
       \nu(Q)^{2/3}\|V\|_{L^6(Q)}^2.
\end{equation}
The cubes form a fixed half-open grid; taking suprema on their closures
does not change any estimate.  The quantity
$\mathcal T_\nu(V)^{1/2}$ is a seminorm, namely the $\ell^2$ norm of
$\nu(Q)^{1/3}\|V\|_{L^6(Q)}$.  This observation will allow us to use
Minkowski's inequality for superpositions of exact extensions.

It suffices first to bound \eqref{eq:sparse-functional} by the right-hand
side of \eqref{eq:src7}, with $\mu$ restricted to a spatial box of side
$O(L^2)$.  Indeed, the three-dimensional unit-cube Sobolev inequality and
H\"older's inequality give, for a fixed finite set of multiindices $a$,
\begin{align}
 \sum_{q\subset Q}\mu(q)^{2/3}\sup_q|U|^2
 &\lesssim
   \sum_a\sum_{q\subset Q}\mu(q)^{2/3}
       \left(\int_q|\partial^aU|^6\right)^{1/3}
       \notag\\
 &\lesssim
   \sum_a\mu(Q)^{2/3}
       \left(\int_Q|\partial^aU|^6\right)^{1/3}.
       \label{eq:sparse-sobolev}
\end{align}
The dyadic $K^2$-grid is nested with the unit grid, and the Sobolev
inequality on a cube controls the supremum on its closure, so no
boundary enlargement is needed in this calculation.  Differentiating an atom
in \eqref{eq:src5} produces only bounded carrier-frequency factors and
native profile derivatives multiplied by bounded powers of $L^{-1}$.
It therefore gives another finite sum of atoms in the same class, with
the same labels and comparable coefficient energy.  Applying the
$\mathcal T$ estimate to these differentiated arrays proves the desired
unit-cube statement.

For this reduction, first localize the original compact packets to
spatial boxes of side $L^2$.  On the stipulated time interval, one such
packet meets only boundedly many enlarged boxes, since its velocity is
bounded.  Thus the energies of the retained arrays sum with bounded
overlap.  This justifies summing the spatially localized estimates and
also the treatment of boundary cubes in \eqref{eq:sparse-sobolev}.

The compact profiles are convenient for this first localization but are
not themselves exact Schr\"odinger extensions.  The following construction
passes to exact fields once, before the iteration.

\begin{lemma}[Exact Fourier layers]\label{lem:sparse-layers}
Fix a bounded native time interval $|\sigma|\le C_{\rm nat}$, any finite
number of profile seminorms on that interval, and any prescribed finite
spatial decay orders there.  After requiring sufficiently many derivatives of the
compact profiles, an array of atoms in \eqref{eq:src5} is a sum-integral
of exact extension arrays, indexed by $k\in\mathbb Z^2$ and $u\in\R$,
times a common phase $e^{iu t/L^2}$.  There are nonnegative common weights
$w(k,u)$, with
\[
 \sum_k\int_\R w(k,u)\,du<\infty,
\]
such that division by $w(k,u)$ gives all the stipulated uniform
seminorm and decay bounds.  For every subarray $A$ the frequency input
$v_{A,k,u}$ of its layer satisfies
\begin{equation}\label{eq:sparse-layer-energy}
 \|v_{A,k,u}\|_2^2
    \le C w(k,u)^2\sum_{\alpha\in A}|c_\alpha|^2.
\end{equation}
Its exact frequencies lie within $O(L^{-1})$ of
$\omega_\alpha+k/L$ for the corresponding labels.
\end{lemma}

\begin{proof}
Use the Fourier transform in all native variables $(y,\sigma)\in\R^3$,
and choose a smooth partition of unity
$\sum_{k\in\mathbb Z^2}\rho(\eta-k)=1$, with $\rho$ supported in
$[-2,2]^2$.  With
\[
 a_{\alpha,k,u}(\eta)
   =\rho(\eta-k)\widehat H_\alpha(\eta,u-|\eta|^2),
\]
Fourier inversion followed by $u=\zeta+|\eta|^2$ gives the exact identity
\begin{equation}\label{eq:sparse-native-layers}
 H_\alpha(y,\sigma)
   =(2\pi)^{-3}\sum_k\int_\R e^{iu\sigma}
       \left[\int_{\R^2}
          e^{i(y\cdot\eta-\sigma|\eta|^2)}
          a_{\alpha,k,u}(\eta)\,d\eta\right]du.
\end{equation}
The carrier phase combines with the inner phase according to
\begin{align}
 x\cdot\omega_\alpha-t|\omega_\alpha|^2
  +\frac{x-b_\alpha-2\omega_\alpha t}{L}\cdot\eta
  -\frac{t}{L^2}|\eta|^2
   & =x\cdot(\omega_\alpha+\eta/L)
       -t|\omega_\alpha+\eta/L|^2
       -b_\alpha\cdot\eta/L.
       \label{eq:sparse-exact-phase}
\end{align}
Thus, with $Ev(x,t)=\int e^{i(x\cdot\xi-t|\xi|^2)}v(\xi)\,d\xi$,
the exact datum for one packet in a fixed layer is
\begin{equation}\label{eq:sparse-layer-datum}
 v_{\alpha,k,u}(\xi)
   = L\,a_{\alpha,k,u}\bigl(L(\xi-\omega_\alpha)\bigr)
         e^{-ib_\alpha\cdot(\xi-\omega_\alpha)}.
\end{equation}
The factor $L$ makes its $L^2$ norm equal to
$\|a_{\alpha,k,u}\|_2$; the inverse Jacobian $L^{-2}$ then produces
the atom's prefactor $L^{-1}$ in the extension field.

Compact support and sufficiently many derivatives of $H_\alpha$ give
any prescribed finite collection of rapidly decaying bounds for its
Fourier transform and its derivatives.  Differentiating
$\widehat H_\alpha(\eta,u-|\eta|^2)$ introduces only polynomial factors
in $\eta$.  On $\eta\in k+[-2,2]^2$ these, including the case
$|u|\asymp |k|^2$, are absorbed by taking more original derivatives.
Explicitly, on this support,
\[
 (1+|k|)^P(1+|u|)^P
 \le C_P\bigl(1+|\eta|+|u-|\eta|^2|\bigr)^{3P}.
\]
After fixing the desired seminorms, the spatial decay orders, and
$C_{\rm nat}$, differentiation and integration by parts introduce only
an additional factor $(1+|k|)^B$ for a fixed finite $B$.  Taking the
original Fourier decay order larger than $3P+B$ therefore permits,
for any subsequently specified large finite $P$, the common weight
\[
 w(k,u)=C_P(1+|k|)^{-P}(1+|u|)^{-P},
\]
so that the desired derivatives of $a_{\alpha,k,u}/w(k,u)$ are uniformly
bounded.  On the fixed bounded native time interval, integration by
parts in $\eta$ also gives the required finite-order spatial decay for
the bracket in \eqref{eq:sparse-native-layers}, after division by the
same weight.  Derivatives of $e^{-i\sigma|\eta|^2}$ cost polynomial
powers of $k$, already included in $B$.  This also absorbs the
displacement $y\simeq2\sigma k$ of a layer's own propagated center:
spatial decay is asserted after normalization by the weight chosen
for the prescribed decay order.  In particular these
exact packet fields decay about the original label trajectories to any
prescribed finite order, uniformly in the layer after normalization.

For completeness, the energy bound uses the full phase-space separation,
not bounded multiplicity of frequencies alone.  In the Gram matrix of
\eqref{eq:sparse-layer-datum}, the entry indexed by $\alpha,\widetilde\alpha$
vanishes unless
$L|\omega_\alpha-\omega_{\widetilde\alpha}|\le C$, since the same
$k$ occurs in both frequency cutoffs.  On an overlap, integration by
parts in the rescaled frequency variable gives, for any fixed large $N$,
\begin{equation}\label{eq:sparse-gram}
 |\langle v_{\alpha,k,u},v_{\widetilde\alpha,k,u}\rangle|
 \le C_N w(k,u)^2
    \mathbf1_{\{L|\omega_\alpha-\omega_{\widetilde\alpha}|\le C\}}
    \left(1+\frac{|b_\alpha-b_{\widetilde\alpha}|}{L}\right)^{-N}.
\end{equation}
The bounded number of labels in every unit phase-space ball makes each
row and column summable.  Schur's test proves
\eqref{eq:sparse-layer-energy}, for an arbitrary subarray as well as
for the entire array.
\end{proof}

We may therefore prove the $\mathcal T$ estimate for the normalized exact
arrays of Lemma~\ref{lem:sparse-layers}.  The phase $e^{iu t/L^2}$ is
common to an entire layer and has modulus one, so it does not change
$\mathcal T$ at the root.  Minkowski's inequality and the summability of
$w$ return the estimate for the original array.  If derivatives are
needed in \eqref{eq:sparse-sobolev}, differentiate the compact atoms
\emph{before} this reduction; no differentiation of an exterior layer
phase inside the iteration is needed.

Fix a very small $\gamma>0$ later.  Layers with $|k|>L^\gamma$ can be
discarded with an arbitrarily small negative-power error, by increasing
$P$ and the number of original derivatives.  To see uniformity, normalize
$m=\|c\|_{\ell^2}=1$ and use the spatial restriction already made.
The number of relevant cells, their total mass, and a trivial bound for
the functional are polynomial in $L$; \eqref{eq:sparse-layer-energy}
and the bounded area of each layer's frequency support give a uniform
field supremum bound after normalization.  The tail sum of $w$ beats
any of these fixed polynomial factors.  Homogeneity restores general
$m,c$.  Henceforth fix one retained layer and suppress $k,u,w$ from
the notation.

\subsection{The iteration and broad saving}

We first specify the scales at which the next estimates will be used.
A child retains one frequency square of side $K^{-1}$ and applies
the cap rescaling in \eqref{eq:prelim-atom-rescaling}.  At depth $j$
set $h=K^j$ and $s=L/h$, and divide out the accumulated field factor
$h^{-1}$.  The transformed labels occupy a fixed bounded square.
The native variables, native time bounds, and uniform profile bounds
are unchanged.

Stop at the last scale before the next division by $K$ would go below
$D^{1/2}$.  Thus every scale satisfies $s\ge D^{1/2}$, and terminal
scales satisfy
\begin{equation}\label{eq:sparse-terminal-scale}
 D^{1/2}\le s<KD^{1/2}.
\end{equation}
The depth $J$ is bounded by a constant depending on $\beta$.
The broad estimate below is stated independently of this stopping
rule; the terminal estimate will use it together with packet sparsity.

At current packet size $s$, partition the bounded square of label
frequencies into $K^{-1}$-squares $\theta$, and write
$V=\sum_\theta V_\theta$.  Each $V_\theta$ is exact.  Its frequencies
are in its label square, enlarged by $O(s^{-1})$, plus the common shift
$k/s$.  We will always have $s\gg K^2$ and $|k|/s=o(1)$.
For a fixed sufficiently large absolute $A$, call $\theta$ active on
$Q$ if
\begin{equation}\label{eq:sparse-active}
 \|V_\theta\|_{L^6(Q)}
      \ge K^{-A}\|V\|_{L^6(Q)}.
\end{equation}
There are $O(K^2)$ squares.  Thus the inactive sum has norm at most
$C K^{2-A}\|V\|_{L^6(Q)}$, which is absorbed by choosing $A$ once and
for all.

There is a fixed large $C_0$ such that either all active centers lie
within $C_0/K$ of a line, or three active squares have direction
determinant at least $cK^{-2}$ throughout their exact supports.  Indeed,
choose a farthest pair of active centers.  If another center lies outside
the indicated strip about their line, the resulting triangle has area
at least a constant times $K^{-2}$.  More precisely its height is at
least $C_0/K$, and its base is at least as long as this height.  Changing
each vertex by $O(K^{-1})$ changes its doubled area by at most a fixed
constant times the base length divided by $K$; large $C_0$ preserves a
fixed portion of the area.  The identity between triangle area and the
determinant of $(2\xi_i,1)$, and the bounded frequency range, give the
claim also for normalized directions.  A common frequency translation
does not change the unnormalized determinant.  Call the two alternatives
\emph{narrow} and \emph{broad}, respectively.

\begin{lemma}[Broad saving]\label{lem:sparse-broad}
There is an absolute $c_b>0$ such that, whenever $K\le s^{c_b}$, the
broad part of the functional in a ball of radius $O(s^2)$ satisfies
\begin{equation}\label{eq:sparse-broad}
 \sum_{Q\text{ broad}}\nu(Q)^{2/3}\|V\|_{L^6(Q)}^2
       \lesssim s^{-c_b}m_0^{2/3}s^{4/3}\mathcal E.
\end{equation}
Here $\nu(B(z,r))\le m_0r^2$ for $r\ge1$, and the exact frequency
data for $V$ and each $V_\theta$ have $L^2$ norms at most
$C\mathcal E^{1/2}$.  The implicit constant can depend on this fixed
$C$ but $c_b$ does not depend on packet differentiability.
\end{lemma}

\begin{proof}
Suppose no positive $c_b$ works.  After normalizing $m_0=\mathcal E=1$
and dividing all fields by a fixed bound for their input norms, this
gives a sequence $s\to\infty$, with $K=s^{o(1)}$, for which the broad
sum is at least $s^{4/3-o(1)}$.  Fixed constants in the radius, support
box and norm bounds are harmless.  There are polynomially many cubes,
and the fields have bounded supremum by Cauchy--Schwarz in frequency.
We may discard polynomially tiny mass and field levels, and then
pigeonhole one active triple and dyadic values
\[
 \nu(Q)\asymp\chi,\qquad \|V\|_{L^6(Q)}\asymp b
\]
on $N$ cubes.  The number of triples and all factors lost here are
$K^{O(1)}$ or logarithmic, hence $s^{o(1)}$.  We obtain
\begin{equation}\label{eq:sparse-saturation-bounds}
 \chi^{2/3}b^2N\ge s^{4/3-o(1)},\qquad
 \chi\le s^{o(1)},\qquad
 \chi N\lesssim s^4.
\end{equation}
The last two statements follow respectively from the ball bound on a
$K^2$-cube and on the ambient ball.

We explain carefully the fourth inequality, furnished by multilinear
restriction.  Each chosen factor has $L^6(Q)$ norm at least $K^{-A}b$,
so some unit cube in $Q$ has supremum at least $K^{-A-1}b$: here
$|Q|^{1/6}=K$.  For each of the three factors, choose one such unit
cube.  Pigeonholing their three relative integer positions within $Q$
leaves a common triple of offsets, with cost $K^{O(1)}$.  Translate
the three fields separately by these fixed offsets.  They now all have
supremum at least $b\,s^{-o(1)}$ on the same unit cube attached to each
surviving $Q$.  These attached cubes are disjoint.  Translation only
modulates the frequency inputs and preserves both their norms and
their supports.

The separate-supremum version of
Lemma~\ref{lem:transverse-product} therefore yields
\begin{equation}\label{eq:sparse-product-count}
 b^3N\le s^{o(1)}.
\end{equation}
Equivalently, this use of the lemma follows by applying a rapidly
decreasing reproducing kernel to each factor separately, integrating
the product on each attached cube, and averaging the trilinear integral
bound over the three independent shifts.  For every fixed triple of
shifts the integration domain is the same and the inputs are exact
with unchanged supports and norms.  The determinant lower bound is
$K^{-2}=s^{-o(1)}$, precisely the regime of that lemma.  Thus an
unspecified dependence of a multilinear constant on transversality is
not being used here.

The four inequalities force saturation of all three sizes.  For
clarity, write $(x,y,n)$ for limiting base-$s$ logarithms of
$(\chi,b,N)$, passing to a subsequence if needed.  The retained levels
and the first inequality give polynomial bounds, so such subsequences
exist.  Equations \eqref{eq:sparse-saturation-bounds} and
\eqref{eq:sparse-product-count} imply
\[
 x\le0,\qquad x+n\le4,\qquad 3y+n\le0,
 \qquad 2x+6y+3n\ge4.
\]
Consequently
\[
 4\le2x+6y+3n\le2x+n\le4+x\le4.
\]
Equality holds throughout, giving $x=0$, $n=4$, and $y=-4/3$.
In the original sequence this says
\begin{equation}\label{eq:sparse-saturated}
 \chi=s^{o(1)},\qquad N=s^{4+o(1)},\qquad b=s^{-4/3+o(1)}.
\end{equation}
In particular the needed lower bounds are
$\chi\ge s^{-o(1)}$, $N\ge s^{4-o(1)}$ and
$b\ge s^{-4/3-o(1)}$.

The attached unit cubes also have the required ball counting bound.
If their locations lie in a radius-$r$ ball, their associated full
$K^2$-cubes lie in a ball of radius $r+CK^2$.  Each associated cube has
mass at least $s^{-o(1)}$, and they are disjoint.  For $r\ge1$ their
number is therefore at most
\[
 s^{o(1)}(r+CK^2)^2\le s^{o(1)}r^2.
\]
The three translated exact fields, the cubes in
\eqref{eq:sparse-saturated}, and the determinant bound thus contradict
Proposition~\ref{prop:transverse}.  A fixed enlargement of the ambient
ball can be removed by replacing $s$ by a fixed multiple and absorbing
the constants into the subpower factors.

This contradiction supplies an absolute positive $c_b$, reducing it
if necessary so that the same exponent appears in both its hypothesis
and its conclusion.  The argument used only the maximum of the total
input norm and the individual subfield input norms.  Normalizing by
that maximum preserves activity ratios.  Thus it is uniform for exact
fields and its exponent is independent of the high decay orders used
to localize packet arrays later.
\end{proof}

\subsection{Narrow decoupling and the child measure}

We record the narrow estimate with its measure transformation, since
the latter is what allows sparsity to survive the iteration.  For any
fixed sufficiently large decay order $N_0$, write
\[
 w_Q(z)=\left(1+\frac{|z-z_Q|}{K^2}\right)^{-N_0}.
\]
On a narrow cube, parabola decoupling gives
\begin{equation}\label{eq:sparse-narrow-decoupling}
 \|V\|_{L^6(Q)}^2
    \lesssim_\epsilon K^\epsilon
          \sum_\theta\|V_\theta\|_{L^6(w_Q)}^2.
\end{equation}
Here the sum can include all children; only active children are needed
in deriving the inequality.

To check the exact decoupling geometry, rotate frequency coordinates
so that the active strip is $|\xi_2-v_2|\le C/K$, and put
$\eta_2=\xi_2-v_2$.  A bounded physical shear
$y_2=x_2-2v_2t$ and a modulation remove the constant and linear normal
terms.  At each fixed $y_2$, the resulting function of $(x_1,t)$ has
Fourier support in
\[
 \{(\xi_1,\tau): |\tau+\xi_1^2|\le C K^{-2}\},
\]
because the remaining normal contribution is $-\eta_2^2$.
Its tangential cap scale is $K^{-1}$.  The $K^{-1}$ grid of original
patch centers has boundedly many centers in each such tangential bin
inside a strip of width $C/K$: the corresponding planar region has
area $O(K^{-2})$.  The fixed enlargement $O(s^{-1})$ preserves this
bounded overlap.  Also, each original patch meets only boundedly many
tangential bins.

For localization, multiply by a Schwartz cutoff $\psi_Q$ bounded below on $Q$
whose Fourier transform is supported in a ball of radius $O(K^{-2})$.
After the bounded shear and rotation, its restriction to each $y_2$
slice has the same order of tangential/time Fourier thickness.  The
thick-support hypotheses of the parabola theorem therefore remain
valid.  Write $F_\theta=\psi_QV_\theta$ in these coordinates, and
$F=\sum_{\theta\text{ active}}F_\theta$.  On each fixed $y_2$ slice,
use the smooth tangential cap projections $P_I$ at interval length
$K^{-1}$ in the usual smooth-cap form of Bourgain--Demeter decoupling
\cite[Theorem~1.1]{BD15}.  Their multipliers depend only on $\xi_1$;
their one-dimensional convolution kernels have uniformly bounded
$L^1$ norms.  Thus, with $\delta\asymp K^{-2}$,
\[
 \|F(\cdot,y_2,\cdot)\|_{L^6_{x_1,t}}^2
 \lesssim_\epsilon K^\epsilon
       \sum_I\|P_IF(\cdot,y_2,\cdot)\|_{L^6_{x_1,t}}^2.
\]
This decouples frequency projections of the localized sum.  To
recover the original, possibly overlapping subfields, let
$\mathcal N(I)$ consist of the patches whose enlarged tangential
supports meet $I$.  The preceding geometry gives
$|\mathcal N(I)|\le C$, and each patch belongs to at most $C$ such
sets; the cutoff's extra $O(K^{-2})$ thickness does not change these
bounds.  Since $P_IF=\sum_{\theta\in\mathcal N(I)}P_IF_\theta$,
the triangle inequality and the uniform convolution bound give
\[
 \sum_I\|P_IF(\cdot,y_2,\cdot)\|_6^2
 \le C\sum_I\sum_{\theta\in\mathcal N(I)}
                 \|F_\theta(\cdot,y_2,\cdot)\|_6^2
 \le C'\sum_\theta\|F_\theta(\cdot,y_2,\cdot)\|_6^2.
\]
In particular no lower bound on any component of a possibly cancelling
sum is required.  Taking the $L^6$ norm in $y_2$ and using Minkowski
now gives
\[
 \left\|\left(\sum_\theta
       \|F_\theta(\cdot,y_2,\cdot)\|_{L^6_{x_1,t}}^2
             \right)^{1/2}\right\|_{L^6_{y_2}}^2
 \le\sum_\theta\|F_\theta\|_{L^6}^2
 \lesssim\sum_\theta\|V_\theta\|_{L^6(w_Q)}^2.
\]
Here $|\psi_Q|^6\lesssim w_Q$; all projections above were estimated
in global slice norms after localization, so no inverse of the
physical cutoff is needed.  Choosing the cutoff with
arbitrarily high fixed decay gives any required $N_0$.  Finally the
inactive sum is absorbed using \eqref{eq:sparse-active}.  This proves
\eqref{eq:sparse-narrow-decoupling} for the stated subfields.

Fix a small padding exponent $\eta>0$.  The weighted norms can be
truncated to $L^\eta Q$, with an arbitrarily small negative-power error,
by taking $N_0$ sufficiently large.  Exact fields here have bounded
frequencies and controlled input norms, hence polynomial global
supremum bounds.  Their measures and domains have polynomial size in
$L$, so this truncation is uniform after energy and mass normalization.

Consider a child square centered at $v$ and its parabolic map
\begin{equation}\label{eq:sparse-child-map}
 \Phi_v(x,t)=\left(\frac{x-2vt}{K},\frac{t}{K^2}\right),
 \qquad s'=s/K,\quad b_\alpha'=b_\alpha/K,\quad
       \omega_\alpha'=K(\omega_\alpha-v).
\end{equation}
After removing a modulation, $V_\theta=K^{-1}V_\theta'\circ\Phi_v$.
This transformation leaves the native packet variables and the
phase-space counting condition unchanged.  If
$\nu'=(\Phi_v)_\#\nu$, then
\begin{equation}\label{eq:sparse-direct-transport}
 \nu'(B(z,r))\le C K^3m_0r^2\quad(r\ge1),
 \qquad
 \int\left(1+\frac{|x'-b_\alpha'-2\omega_\alpha't'|}{s'}\right)^{-M}
            d\nu'
 \le \frac{m_0s^3}{D}.
\end{equation}
Indeed, the inverse image of a radius-$r$ ball has dimensions
$O(Kr),O(Kr),O(K^2r)$ and is covered by $O(K)$ balls of radius $O(Kr)$.
This proves the first bound.  The normalized transverse distance in
the second integral equals its parent value pointwise, proving the
second bound.  Its right-hand side is exactly
$(K^3m_0)(s')^3/D$.

The measure directly pushed forward is slightly modified to account
for the weighted norms.  Group parent cubes $Q$ by the child
$K^2$-cube $P$ containing $\Phi_v(z_Q)$.  Their total mass is at most
$\nu'(CP)$, since their images lie in $CP$.  The sets
$\Phi_v(L^\eta Q)$ lie in $CL^\eta P$ and have overlap at most
$CL^{3\eta}$.  The Jacobian and field normalization give
\begin{equation}\label{eq:sparse-jacobian}
 \left(\int_{L^\eta Q}|V_\theta|^6\right)^{1/3}
   =K^{-2/3}
       \left(\int_{\Phi_v(L^\eta Q)}|V_\theta'|^6\right)^{1/3},
\end{equation}
because $dx\,dt=K^4dx'\,dt'$ and $|V_\theta|^6=K^{-6}|V_\theta'|^6$.
H\"older with exponents $3/2$ and $3$, first within each group, yields
\begin{align}
 &\sum_Q\nu(Q)^{2/3}
      \left(\int_{L^\eta Q}|V_\theta|^6\right)^{1/3}
       \notag\\
 &\hspace{1cm}\le C L^{O(\eta)}K^{-2/3}
       \sum_P\nu'(CP)^{2/3}
          \left(\int_{CL^\eta P}|V_\theta'|^6\right)^{1/3}.
          \label{eq:sparse-group-holder}
\end{align}
There is no factor for the number of parent cubes in a group.

To put the last sum on the child grid, cover $CP$ by child cubes $P+a$
with $a$ in a fixed finite set of lattice shifts, and cover
$CL^\eta P$ by $P+b$ with $b$ in a set of $O(L^{3\eta})$ shifts.
Both shift lengths are at most $CL^\eta K^2$.  Subadditivity of the
powers $2/3$ and $1/3$ bounds a summand by the sum of
\[
 \nu'(P+a)^{2/3}
        \left(\int_{P+b}|V_\theta'|^6\right)^{1/3}.
\]
For a vector $d$, let $(T_d)_\#\nu'$ denote translation of the measure
by $d$, and define
\begin{equation}\label{eq:sparse-child-measure}
 \nu_{\theta}
    =\sum_{d\in\mathcal D}(T_d)_\#\nu',
 \qquad
 \mathcal D=\{b-a\text{ for the shifts above}\}.
\end{equation}
For $d=b-a$, the mass of $P+b$ under $(T_d)_\#\nu'$ equals
$\nu'(P+a)$.  Summing with the possible multiplicities, which are
$L^{O(\eta)}$, transforms \eqref{eq:sparse-group-holder} into
\begin{equation}\label{eq:sparse-child-functional}
 \sum_Q\nu(Q)^{2/3}
      \left(\int_{L^\eta Q}|V_\theta|^6\right)^{1/3}
 \le C L^{O(\eta)}K^{-2/3}\mathcal T_{\nu_\theta}(V_\theta').
\end{equation}

The number of translates is $L^{O(\eta)}$, so their ball bounds hold
with parameter
\begin{equation}\label{eq:sparse-child-density}
 m_1=C K^3L^{O(\eta)}m_0.
\end{equation}
They also preserve the sparsity bound with this same enlarged
parameter.  In fact, for a shift $d=(d_x,d_t)$, the child's transverse
displacement changes by $d_x-2\omega_\alpha'd_t$.  Child velocities
are bounded, and we impose
\begin{equation}\label{eq:sparse-shift-smallness}
 L^\eta K^2\ll s'
\end{equation}
at every existing child scale.  The weight
$(1+|x'-b_\alpha'-2\omega_\alpha't'|/s')^{-M}$ then changes by at most
a fixed factor under any translation in \eqref{eq:sparse-child-measure}.
Summing \eqref{eq:sparse-direct-transport} over these shifts proves
\[
 \int\left(1+\frac{|x'-b_\alpha'-2\omega_\alpha't'|}{s'}\right)^{-M}
       d\nu_\theta
 \le m_1(s')^3/D.
\]
Their time shifts are also $o(s')$, so a time range $O((s')^2)$ remains
of that order.  Combining \eqref{eq:sparse-narrow-decoupling} and
\eqref{eq:sparse-child-functional} proves the narrow recursion, with
factor $C K^{-2/3+\epsilon}L^{O(\eta)}$ per generation.  Its principal
scale factors satisfy the exact cancellation
\begin{equation}\label{eq:sparse-scale-cancellation}
 K^{-2/3}(K^3m_0)^{2/3}(s/K)^{4/3}
       =m_0^{2/3}s^{4/3}.
\end{equation}

\subsection{Terminal estimate and summation}

\begin{proof}[Proof of Proposition~\ref{prop:sparse}]
We apply the preceding estimates on the scales just specified.
At any node the preceding transport calculation gives a measure $\nu$
with time range $O(s^2)$, ball parameter $m_j$, and packet moment bound
$m_js^3/D$, where
\begin{equation}\label{eq:sparse-node-parameters}
 m_j\le C^j h^3m L^{Cj\eta}.
\end{equation}
All support boxes and cell counts are polynomial in $L$.  We may apply
the estimates at radius $O(s^2)$ spatial box by spatial box.  Here exact
packets replace the compact packets used at the root: retain only
packets whose trajectories meet a fixed enlargement of the box, and
discard the others using their prescribed rapid decay.  On the time
range $O(s^2)$, each retained packet is used in boundedly many boxes.
Its coefficient energy therefore has bounded overlap, also for every
child subarray.  A discarded packet has native transverse distance
$\gtrsim s$ on the box.  Cauchy--Schwarz, the phase-space counting
condition and a sufficiently high decay order make the total error an
arbitrarily small negative power of $L$, since
$s\ge L^{\tau_0/2}$.  This is the localization mechanism of
Lemma~\ref{lem:packets}, applied to the exact arrays constructed above.

Whenever this localization is performed, use the same retained labels
in the total field and in all its subfields.  Negligible field-level
cubes are absorbed in the additive error; on the remaining cubes the
activity tests tolerate that error, or can be made afresh for the
localized exact array.  Lemma~\ref{lem:sparse-broad} then uses the
maximum of its total and subfield input norms, bounded by a fixed
multiple of the retained coefficient norm by
\eqref{eq:sparse-layer-energy}.  In particular increasing the later
decay orders cannot reduce the absolute broad saving.

At every nonterminal node the broad contribution is bounded by
\eqref{eq:sparse-broad}, and the narrow contribution is passed to
children by \eqref{eq:sparse-child-functional}.  Child arrays partition
the parent labels.  Hence their squared coefficient norms sum to the
parent energy; there is no loss from the number of children.

It remains to estimate a terminal node.  By rapid decay keep only
$K^2$-cubes lying within transverse distance $L^\eta s$ of at least
one label trajectory.  Taking the decay order sufficiently high
again makes the omitted contribution negligible.  Restrict $\nu$ to
these retained cubes.  If a ball $B$ of radius $s$ meets the retained
support, choose one retained cube it meets and a packet near that
cube.  Since $K^2\ll s$ and velocities are bounded, every point of
$B$ is at normalized transverse distance $O(L^\eta)$ from that one
packet.  Its moment bound therefore gives
\begin{equation}\label{eq:sparse-terminal-mass}
 \nu(B)\le C L^{CM\eta}\frac{m_js^3}{D}.
\end{equation}
The same estimate holds for fixed enlargements of $B$.  No sum over
packets occurs: one packet suffices for each ball.

Normalize the node coefficient energy to one and split retained cubes
into dyadic classes
$\nu(Q)\asymp m_j\chi$.  The upper range is
$\chi\lesssim K^4$ by the ball bound.  Polynomially tiny levels are
negligible: there are polynomially many tested cells, and a uniform
polynomial supremum bound follows from the exact input energy.
Thus only $O(\log L)$ levels need be kept.  In each $Q$ choose a unit
cube $q_Q$ with the largest field supremum.  Then
\begin{equation}\label{eq:sparse-terminal-unit-reduction}
 \sum_{Q\text{ in the class}}\nu(Q)^{2/3}\|V\|_{L^6(Q)}^2
 \lesssim K^2(m_j\chi)^{2/3}
       \sum_Q\sup_{q_Q}|V|^2.
\end{equation}
The selected unit cubes need not themselves carry measure.  Their
counts are instead controlled by their disjoint full $Q$-cubes, each
of mass comparable to $m_j\chi$.  For $r\ge1$ this gives
\[
 \#\{q_Q:q_Q\cap B(z,r)\ne\varnothing\}
 \lesssim \chi^{-1}(r+CK^2)^2
 \lesssim K^4\chi^{-1}r^2.
\]
At radius $s$, using \eqref{eq:sparse-terminal-mass} and $K^2\ll s$
instead gives occupancy at most
$C L^{CM\eta}\chi^{-1}s^3/D$.
Thus the parameters in Lemma~\ref{lem:fractal}, with ambient radius
$S\asymp s^2$, can be taken as
\[
 \gamma_0\lesssim K^4\chi^{-1},\qquad
 \lambda_0\lesssim L^{CM\eta}\chi^{-1}s^3/D.
\]
Its unit-cube supremum version and
\eqref{eq:sparse-terminal-unit-reduction} yield
\begin{align}
 \mathcal T_\nu(V)
 &\lesssim_\epsilon
    L^{CM\eta}K^{10/3}L^\epsilon
       m_j^{2/3}s^{5/3}D^{-1/3}\|c_{\rm node}\|_{\ell^2}^2
       \notag\\
 &=L^{CM\eta}K^{10/3}O_\epsilon(L^\epsilon)
       m_j^{2/3}s^{4/3}(s/D)^{1/3}
       \|c_{\rm node}\|_{\ell^2}^2.
       \label{eq:sparse-terminal}
\end{align}
Here the logarithmic number of weight classes is absorbed in
$L^\epsilon$.  The cancellation of the weight parameter is complete:
\[
 \chi^{2/3}(\chi^{-1}\chi^{-1})^{1/3}=1.
\]
The displayed $K^{10/3}$ consists of $K^2$ in
\eqref{eq:sparse-terminal-unit-reduction} and $K^{4/3}$ from the
counting parameter.  Allowing additional fixed cube or grid
enlargements merely replaces it by $K^{C_1}$ for an absolute $C_1$.
Spatial localization is as above, so this argument applies even when
the node's spatial support extends beyond a single $s^2$-box.

We finish with an explicit order of choices.  Bounded $L$ causes no
difficulty, and $\tau_0>1$ permits no unbounded admissible scales, so
assume $0<\tau_0\le1$.  Fix first the absolute numbers
\begin{equation}\label{eq:sparse-saving-choice}
 a_* =\min\{c_b/2,1/6\},\qquad c_s=a_*/8.
\end{equation}
At a broad leaf $s\ge D^{1/2}$ gives
$s^{-c_b}\le D^{-c_b/2}$, while
\eqref{eq:sparse-terminal-scale} gives the terminal gain
\[
 (s/D)^{1/3}\le K^{1/3}D^{-1/6}.
\]
Thus both gains are at least $D^{-a_*}$ before the losses just recorded.

Now choose $\beta>0$ sufficiently small in terms of $\tau_0,c_b,a_*$.
It can satisfy, with fixed strict margins,
\[
 \beta<\frac{\tau_0c_b}{4},\qquad
 2\beta<\frac{\tau_0}{4},\qquad
 (C_1+1/3)\beta<\frac{\tau_0a_*}{16}.
\]
For all large $L$, these imply $K\le s^{c_b}$ at every broad node and
allow the terminal $K$ factors to consume at most $D^{a_*/8}$, after
fixed constants are absorbed.  The depth bound $J=O(\beta^{-1})$ is
now fixed.  Choose the small losses in decoupling and the refined
fractal bound so that their total product throughout a path is at
most $D^{a_*/8}$.  Next choose $\eta>0$, depending also on $J$ and
$M$, so that all padding losses together cost at most $D^{a_*/8}$
and
$2\beta+\eta<\tau_0/2$.  This last inequality implies
\eqref{eq:sparse-shift-smallness} at every existing child scale, since
each such scale is at least $D^{1/2}\ge L^{\tau_0/2}$.
Finally choose $\gamma<\tau_0/4$, so that $|k|/s=o(1)$ for every
retained layer, and then choose the frequency-tail exponent, the
physical decay orders, and the original profile differentiability.
The latter choices make every discarded tail as small a negative
power as required; they do not change $c_b$ or $c_s$.

To see that these choices close the entire tree, a depth-$j$ path
contributes, before its broad or terminal gain, at most
\[
 C^j L^{Cj\eta+O(\epsilon)}h^{-2/3},
 \qquad h=K^j,
\]
times its node estimate.  On the other hand
\eqref{eq:sparse-node-parameters} gives
\[
 m_j^{2/3}s^{4/3}
 \le C^jL^{Cj\eta}h^{2/3}m^{2/3}L^{4/3}.
\]
The $h$ powers cancel, as anticipated in
\eqref{eq:sparse-scale-cancellation}.  At a fixed depth the total
branch energy is bounded by a fixed constant to that depth times the
original energy, including the spatial localizations.  Summing over
the bounded depth therefore costs only a constant depending on
$\tau_0$.  The allotted losses leave more than the saving $D^{-c_s}$
in \eqref{eq:sparse-saving-choice}.  Finally, all tail errors can be
made negligible relative to
$m^{2/3}L^{4/3}D^{-c_s}\|c\|_{\ell^2}^2$, even after summing the
tree: after normalization there are only polynomially many supported
cells, boxes and branches, while the tail exponent is chosen last.

This proves the desired bound for $\mathcal T$.  Minkowski in
Lemma~\ref{lem:sparse-layers}, the Sobolev reduction
\eqref{eq:sparse-sobolev}, and the initial spatial energy summation
give \eqref{eq:src7}, with the stated derivative and uniformity
properties.
\end{proof}
